%-----------------------------------------------------------------------------------------------------------------------------
% 1. Spolupráca s vedúcim                 max. 20
% a. Efektivita konzultácií               max. 15
\renewcommand{\bodyKonEfektivita}{} 
% - Do 15b - Študent priebežne konzultoval postup a výsledky práce s vedúcim, riadil sa pokynmi vedúceho a vykonal korekcie ku všetkým pripomienkam  
% - Do 11b - Študent priebežne konzultoval postup a výsledky práce s vedúcim, ale viaceré pokyny alebo pripomienky nerešpektoval
% - Do 7b - Študent priebežne konzultoval postup a výsledky, ale neriadil sa pokynmi vedúceho a výsledky prevažne nezodpovedajú predstave vedúceho
% - 0b - Študent nekonzultoval postup ani výsledky, a tie nezodpovedajú predstave vedúceho


% b. Iniciatíva pri získavaní podkladov   max. 5
\renewcommand{\bodyKonIniciativa}{} 
% - Do 5b - Študent využil okrem odporúčanej literatúry aj iné literárne zdroje a publikácie (knihy, odborné časopisy, vedecké články)
% - Do 3b - Študent využil odporúčanú literatúru a elektronické zdroje
% - Do 1b - Študent nedostatočne využil odporúčanú literatúru a elektronické zdroje

%------------------------------------------------------------------------------------------------------------------------------
% 2. Splnenie cieľov zadania              max. 20
% a. Kvantitatívne plnenie úloh           max. 10
\renewcommand{\bodyKvantita}{}  
% - Do 10b - Všetky úlohy zadania boli riešené v plnom rozsahu
% - Do 5b - Niektorá z úloh si vyžadovala viac priestoru než jej študent venoval

% b. Kvalitatívne plnenie úloh            max. 10
\renewcommand{\bodyKvalita}{}  
% - Do 10b - Výsledky realizovaných úloh úplne spĺňajú očakávania vedúceho
% - Do 7b - Niektorá z úloh, bez ohľadu na rozsah vykonanej práce, nemá výstupy v požadovanej kvalite
% - 0b - Prevažná časť výstupov realizovaných úloh nespĺňa očakávania vedúceho

%------------------------------------------------------------------------------------------------------------------------------
% 3. Kvalita priebežnej práce              max. 30
% a. Organizácia práce                     max. 10
\renewcommand{\bodyPrOrganizacia}{} 
% - Do 10b - Študent na úlohách pracoval pravidelne, rovnomerne si rozložil objem vykonávanej práce. Počas celej doby riešenia práce nikdy nestagnoval a vždy bolo možné pozorovať progres.
% - Do 7b - Študent na úlohách pracoval nepravidelne alebo nárazovo a občas stagnoval. Nemalo to však výrazný dopad na kvalitu vykonávanej práce.
% - Do 4b - Študent na úlohách pracoval nepravidelne alebo nárazovo a výrazne stagnoval, čo sa prejavilo na kvalite vykonávanej práce.

% b. Efektivita práce                      max. 20
\renewcommand{\bodyPrEfektivita}{}
% - Do 20b - Študent efektívne využíval a vhodne prepájal potrebné metódy a nástroje, a teda maximum svojho úsilia mohol vynaložiť na kvalitatívne a kvantitatívne plnenie úloh.
% - Do 15b - Študent venoval väčšinu svojho úsilia na zvládanie potrebných podporných prostriedkov alebo na vykonávanie okrajových činností.
% - Do 10b - Študent venoval väčšinu svojho úsilia na zvládanie okrajových činností alebo na tvorbu podporných systémov. Plnenie podstaty úloh zadania venoval len minimum priestoru.

%------------------------------------------------------------------------------------------------------------------------------
% 4. Dodržanie časového harmonogramu       max. 10
\renewcommand{\bodyHarmonogram}{} 
% - Do 10b - Študent dodržoval dohodnuté termíny, čiastkové výstupy prezentoval vedúcemu podľa dohody a prácu odovzdal načas.
% - Do 6b - Študent prevažne dodržoval dohodnuté termíny. Pracoval s časovým sklzom, no konečnú prácu odovzdal načas.
% - Do 3b - Študent pracoval s výrazným časovým sklzom a prácu odovzdal až počas predĺženého termínu.

%------------------------------------------------------------------------------------------------------------------------------
% 5. Samostatný prístup k práci            max. 20
\renewcommand{\bodySamostatnost}{}
% - Do 20b - Študent bol iniciatívny, sám navrhoval postupy práce a riešenia, ktoré zodpovedali predstave vedúceho. Všetky potrebné poznatky si samostatne doštudoval.
% - Do 12b - Študent iniciatívne pristupoval k riešeniu práce, no nenavrhoval vhodné postupy a riešenia. Spoliehal sa prevažne na pokyny vedúceho.
% - Do 8b - Študent sa pri riešení práce iniciatívne neprejavoval. Pracoval len podľa pokynov vedúceho.
